\documentclass[10pt,a4paper]{amsart}
\usepackage{amsmath,amssymb,amsthm,graphicx}
\usepackage{verbatim}
\newenvironment{reference}{\comment}{\endcomment}
%\newenvironment{reference}{}{}
\newenvironment{slogan}{\comment}{\endcomment}
\newenvironment{history}{\comment}{\endcomment}

% For commutative diagrams we use Xy-pic
\usepackage[all]{xy}

% We use 2cell for 2-commutative diagrams.
\xyoption{2cell}
% Theorem environments.
%
\theoremstyle{plain}
\newtheorem{theorem}[subsection]{Theorem}
\newtheorem{proposition}[subsection]{Proposition}
\newtheorem{lemma}[subsection]{Lemma}

\theoremstyle{definition}
\newtheorem{definition}[subsection]{Definition}
\newtheorem{example}[subsection]{Example}
\newtheorem{exercise}[subsection]{Exercise}
\newtheorem{situation}[subsection]{Situation}

\theoremstyle{remark}
\newtheorem{remark}[subsection]{Remark}
\newtheorem{remarks}[subsection]{Remarks}

\numberwithin{equation}{subsection}

% Macros
%
\def\lim{\mathop{\mathrm{lim}}\nolimits}
\def\colim{\mathop{\mathrm{colim}}\nolimits}
\def\Spec{\mathop{\mathrm{Spec}}}
\def\Hom{\mathop{\mathrm{Hom}}\nolimits}
\def\Ext{\mathop{\mathrm{Ext}}\nolimits}
\def\SheafHom{\mathop{\mathcal{H}\!\mathit{om}}\nolimits}
\def\SheafExt{\mathop{\mathcal{E}\!\mathit{xt}}\nolimits}
\def\Sch{\mathit{Sch}}
\def\Mor{\mathop{\mathrm{Mor}}\nolimits}
\def\Ob{\mathop{\mathrm{Ob}}\nolimits}
\def\Sh{\mathop{\mathit{Sh}}\nolimits}
\def\NL{\mathop{N\!L}\nolimits}
\def\CH{\mathop{\mathrm{CH}}\nolimits}
\def\proetale{{pro\text{-}\acute{e}tale}}
\def\etale{{\acute{e}tale}}
\def\QCoh{\mathit{QCoh}}
\def\Ker{\mathop{\mathrm{Ker}}}
\def\Im{\mathop{\mathrm{Im}}}
\def\Coker{\mathop{\mathrm{Coker}}}
\def\Coim{\mathop{\mathrm{Coim}}}

% Boxtimes
%
\DeclareMathSymbol{\boxtimes}{\mathbin}{AMSa}{"02}

%
% Macros for moduli stacks/spaces
%
\def\QCohstack{\mathcal{QC}\!\mathit{oh}}
\def\Cohstack{\mathcal{C}\!\mathit{oh}}
\def\Spacesstack{\mathcal{S}\!\mathit{paces}}
\def\Quotfunctor{\mathrm{Quot}}
\def\Hilbfunctor{\mathrm{Hilb}}
\def\Curvesstack{\mathcal{C}\!\mathit{urves}}
\def\Polarizedstack{\mathcal{P}\!\mathit{olarized}}
\def\Complexesstack{\mathcal{C}\!\mathit{omplexes}}
% \Pic is the operator that assigns to X its picard group, usage \Pic(X)
% \Picardstack_{X/B} denotes the Picard stack of X over B
% \Picardfunctor_{X/B} denotes the Picard functor of X over B
\def\Pic{\mathop{\mathrm{Pic}}\nolimits}
\def\Picardstack{\mathcal{P}\!\mathit{ic}}
\def\Picardfunctor{\mathrm{Pic}}
\def\Deformationcategory{\mathcal{D}\!\mathit{ef}}
\usepackage{cleveref}
\crefname{originalreference}{original reference}{original references}
\Crefname{originalreference}{Original reference}{Original references}
\makeatletter
\newlabel{lemma-Aut-functor-RS}{{06JU}{source}{}{sourcepointer-1}{}}
\newlabel{lemma-Aut-functor-RS@cref}{{[originalreference][0][]06JU}{[1][0][]source}{}{}{}}
\makeatother
\crefname{originalreference}{original reference}{original references}
\Crefname{originalreference}{Original reference}{Original references}
\makeatletter
\newlabel{lemma-localize-RS}{{06JC}{source}{}{sourcepointer-2}{}}
\newlabel{lemma-localize-RS@cref}{{[originalreference][0][]06JC}{[1][0][]source}{}{}{}}
\makeatother
\crefname{originalreference}{original reference}{original references}
\Crefname{originalreference}{Original reference}{Original references}
\makeatletter
\newlabel{lemma-tangent-space-functor}{{06IA}{source}{}{sourcepointer-3}{}}
\newlabel{lemma-tangent-space-functor@cref}{{[originalreference][0][]06IA}{[1][0][]source}{}{}{}}
\makeatother
\crefname{originalreference}{original reference}{original references}
\Crefname{originalreference}{Original reference}{Original references}
\makeatletter
\newlabel{remark-localize-cofibered-groupoid}{{06GU}{source}{}{sourcepointer-4}{}}
\newlabel{remark-localize-cofibered-groupoid@cref}{{[originalreference][0][]06GU}{[1][0][]source}{}{}{}}
\makeatother
\title{The tangent space of the automorphism group of an object is a k-vector space whose addition is composition}
\author{The Stacks Project Authors}
\date{stacks-project:formal-defos.tex:local-snapshot-20261012; GNU FDL 1.2 or later}
\begin{document}
\maketitle
\noindent\textit{Editorial note.} All delivered mathematical text is the author's own, copied from the pinned Stacks Project chapter file formal-defos.tex; no statement or proof step was rewritten, reordered or removed.\par\medskip
\noindent\textit{Editorial note.} The counted proof is the author's complete argument: the localization of the automorphism functor to a deformation functor, the identification of addition with composition through the (RS) decomposition and a group-theoretic argument, and the comparison of the tangent spaces of the automorphism functors by unwinding the definitions.\par\medskip
\noindent\textit{Editorial note.} The label lemma-Aut-functor-tangent-space names the counted statement, and the proof environment immediately following it in the source is delivered verbatim as its proof.\par\medskip
\noindent\textit{Editorial note.} The counted statement is carried by the author's own lemma environment in the chapter source.\par\medskip
\noindent\textit{Editorial note.} Pre-open gate: python operations/agent-authoring/tools/precard\_gate.py --source formal-defos.tex --label lemma-Aut-functor-tangent-space --start 5928 --end 5998 --local-id 28015 returned RESULT=OK (L28015 free, occupancy FREE, public interval clear). Candidate taken from sources/stacks-formal-defos/candidates-halfA.json.\par\medskip
\noindent\textit{Editorial note.} The statement and proof are copied from the pinned chapter file; the cited inputs are kept as external pointers with their canonical Stacks tags.\par\medskip
\noindent\textit{Allowed standard prerequisites (editorial).}\begin{itemize}
\item Localization of a functor at an element. Same chapter, the remark localizing a cofibred groupoid along an element of the fibre over k, producing a predeformation functor. Supplies the predeformation functor F = Aut(x) localized at the identity of x\_0.
\item Deformation functor structure of the localized automorphism functor. Same chapter, the lemmas that the automorphism functor satisfies (RS) and that localization preserves (RS). Supply the fact that the localized functor is a deformation functor, so that (RS) applies to it.
\item Vector space structure on the tangent space of a deformation functor. Same chapter, the lemma providing the natural k-vector space structure on the tangent space of a deformation functor. Supplies the vector space structure on F(k[epsilon]) whose addition is then identified with composition.
\end{itemize}\medskip
\ifdefined\hypertarget\hypertarget{sourcepointer-1}{}\fi
\noindent\textit{Source reference.} Original reference 06JU: sources/stacks-formal-defos/pinned-original-source.tex, source lines 5913--5913. Prerequisite pointer invoked at proof line(s) 5952; it names a standard fact of the same chapter that the counted argument uses, and every step of the counted proof built on it is retained verbatim in the delivered body.\par\medskip
\ifdefined\hypertarget\hypertarget{sourcepointer-2}{}\fi
\noindent\textit{Source reference.} Original reference 06JC: sources/stacks-formal-defos/pinned-original-source.tex, source lines 5349--5349. Prerequisite pointer invoked at proof line(s) 5952; it names a standard fact of the same chapter that the counted argument uses, and every step of the counted proof built on it is retained verbatim in the delivered body.\par\medskip
\ifdefined\hypertarget\hypertarget{sourcepointer-3}{}\fi
\noindent\textit{Source reference.} Original reference 06IA: sources/stacks-formal-defos/pinned-original-source.tex, source lines 3488--3488. Prerequisite pointer invoked at proof line(s) 5957; it names a standard fact of the same chapter that the counted argument uses, and every step of the counted proof built on it is retained verbatim in the delivered body.\par\medskip
\ifdefined\hypertarget\hypertarget{sourcepointer-4}{}\fi
\noindent\textit{Source reference.} Original reference 06GU: sources/stacks-formal-defos/pinned-original-source.tex, source lines 1579--1579. Prerequisite pointer invoked at proof line(s) 5948; it names a standard fact of the same chapter that the counted argument uses, and every step of the counted proof built on it is retained verbatim in the delivered body.\par\medskip
\begin{lemma}
\label{lemma-Aut-functor-tangent-space}
Let $\mathcal{F}$ be a category cofibered in groupoids over
$\mathcal{C}_\Lambda$ satisfying (RS). Let
$x \in \Ob(\mathcal{F}(A))$. Let $x_0$ be a pushforward of $x$ to
$\mathcal{F}(k)$.
\begin{enumerate}
\item $T_{\text{id}_{x_0}} \mathit{Aut}(x)$ has a natural $k$-vector
space structure such that addition agrees with composition in
$T_{\text{id}_{x_0}} \mathit{Aut}(x)$. In particular, composition in
$T_{\text{id}_{x_0}} \mathit{Aut}(x)$ is commutative.
\item There is a canonical isomorphism
$T_{\text{id}_{x_0}} \mathit{Aut}(x) \to
T_{\text{id}_{x_0}} \mathit{Aut}(x_0)$
of $k$-vector spaces.
\end{enumerate}
\end{lemma}

\begin{proof}
We apply
Remark \ref{remark-localize-cofibered-groupoid}
to the functor $\mathit{Aut}(x) : \mathcal{C}_A \to \textit{Sets}$
and the element $\text{id}_{x_0} \in \mathit{Aut}(x)(k)$ to get
a predeformation functor $F = \mathit{Aut}(x)_{\text{id}_{x_0}}$. By
Lemmas \ref{lemma-Aut-functor-RS} and \ref{lemma-localize-RS}
$F$ is a deformation functor. By definition
$T_{\text{id}_{x_0}} \mathit{Aut}(x) = TF = F(k[\epsilon])$
which has a natural
$k$-vector space structure specified by Lemma
\ref{lemma-tangent-space-functor}.

\medskip\noindent
Addition is defined as the composition
$$
F(k[\epsilon]) \times F(k[\epsilon]) \longrightarrow
F(k[\epsilon] \times_k k[\epsilon]) \longrightarrow
F(k[\epsilon])
$$
where the first map is the inverse of the bijection guaranteed by (RS) and the
second is induced by the $k$-algebra map
$k[\epsilon] \times_k k[\epsilon] \to k[\epsilon]$
which maps $(\epsilon, 0)$ and $(0, \epsilon)$ to $\epsilon$.
If $A \to B$ is a ring map in $\mathcal{C}_\Lambda$, then $F(A) \to F(B)$
is a homomorphism where $F(A) = \mathit{Aut}(x)_{\text{id}_{x_0}}(A)$ and
$F(B) = \mathit{Aut}(x)_{\text{id}_{x_0}}(B)$ are groups under
composition. We conclude that
$+ : F(k[\epsilon]) \times F(k[\epsilon])\to F(k[\epsilon])$
is a homomorphism where $F(k[\epsilon])$ is regarded as a
group under composition. With $\text{id} \in F(k[\epsilon])$ the
unit element we see that $+(v, \text{id}) =
+(\text{id}, v) = v$ for any $v \in F(k[\epsilon])$ because
$(\text{id}, v)$ is the pushforward of $v$ along the ring map
$k[\epsilon] \to k[\epsilon] \times_k k[\epsilon]$ with
$\epsilon \mapsto (\epsilon, 0)$.
In general, given a group $G$ with multiplication $\circ$
and $+ : G \times G \to G$ is a homomorphism such that
$+(g, 1) = +(1, g) = g$, where $1$ is the identity of $G$, then $+ = \circ$.
This shows addition in the $k$-vector space structure on $F(k[\epsilon])$
agrees with composition.

\medskip\noindent
Finally, (2) is a matter of unwinding the definitions. Namely
$T_{\text{id}_{x_0}} \mathit{Aut}(x)$ is the set of
automorphisms $\alpha$ of the pushforward of $x$ along
$A \to k \to k[\epsilon]$ which are trivial modulo $\epsilon$.
On the other hand $T_{\text{id}_{x_0}} \mathit{Aut}(x_0)$ is the set of
automorphisms of the pushforward of $x_0$ along
$k \to k[\epsilon]$ which are trivial modulo $\epsilon$.
Since $x_0$ is the pushforward of $x$ along $A \to k$ the result
is clear.
\end{proof}

\end{document}
